% theory.tex
本節基於 Werner Soedel《Vibrations of Shells and Plates》的 Mindlin 板理論框架，建立自由振動的能量泛函，並透過 Ritz 法進行離散。

\subsection{五步推導：從三維彈性體到二維宏觀能量}

\subsubsection{第一步：微觀能量密度}
在彈性體內一點 $(x,y,z)$ 處，
\begin{align}
u_{0b} &= \frac{1}{2} (\sigma_x \varepsilon_x + \sigma_y \varepsilon_y + \tau_{xy} \gamma_{xy}) \label{eq:u0b} \\
u_{0s} &= \frac{1}{2} (\tau_{xz} \gamma_{xz} + \tau_{yz} \gamma_{yz}) \label{eq:u0s} \\
T(t)  &= \frac{1}{2} \int_V \rho \left[ (\partial u/\partial t)^2 + (\partial v/\partial t)^2 + (\partial w/\partial t)^2 \right] dV \label{eq:T0}
\end{align}

\subsubsection{第二步：本構關係與運動學速度}
代入平面應力本構與 Mindlin 位移場 $u = z\psi_x$, $v = z\psi_y$, $w = w$，得：
\begin{align}
u_{0b} &= \frac{1}{2} \left\{ \frac{E}{1-\nu^2} [\varepsilon_x^2 + \varepsilon_y^2 + 2\nu \varepsilon_x \varepsilon_y] + G \gamma_{xy}^2 \right\} \\
u_{0s} &= \frac{1}{2} G (\gamma_{xz}^2 + \gamma_{yz}^2) \\
T(t)  &= \frac{1}{2} \int_\Omega \left\{ \int_{-h/2}^{h/2} \rho \left[ \dot{w}^2 + z^2 (\dot{\psi}_x^2 + \dot{\psi}_y^2) \right] dz \right\} d\Omega
\end{align}

\subsubsection{第三步：幾何變形假設}
面內應變沿厚度線性分佈，橫向剪應變均勻：
\begin{align}
\varepsilon_x &= z \frac{\partial \psi_x}{\partial x},\quad
\varepsilon_y = z \frac{\partial \psi_y}{\partial y},\quad
\gamma_{xy} = z \left( \frac{\partial \psi_x}{\partial y} + \frac{\partial \psi_y}{\partial x} \right) \\
\gamma_{xz} &= \frac{\partial w}{\partial x} + \psi_x,\quad
\gamma_{yz} = \frac{\partial w}{\partial y} + \psi_y
\end{align}
代入後得到應變能密度顯式含 $z$ 與不含 $z$ 的部分。

\subsubsection{第四步：厚度積分與修正}
對 $z$ 積分後定義彎曲剛度 $D$ 與剪切修正係數 $\kappa^2 = 5/6$：
\begin{align}
D &= \frac{E h^3}{12(1-\nu^2)},\quad
D^s = \kappa^2 G h \\
I_2 &= \frac{\rho h^3}{12}
\end{align}

\subsubsection{第五步：宏觀能量泛函}
最後得到彎曲應變能 $U_b$、橫向剪切應變能 $U_s$ 及最大動能 $T_{\max}$（簡諧運動 $w = W \cos\omega t$ 等）：
\begin{align}
U_b &= \frac{1}{2} \int_\Omega \left\{ D \left[ \left(\frac{\partial \Psi_x}{\partial x}\right)^2 + \left(\frac{\partial \Psi_y}{\partial y}\right)^2 + 2\nu \frac{\partial \Psi_x}{\partial x}\frac{\partial \Psi_y}{\partial y} \right] \right. \nonumber \\
&\quad \left. + D\frac{1-\nu}{2} \left( \frac{\partial \Psi_x}{\partial y} + \frac{\partial \Psi_y}{\partial x} \right)^2 \right\} d\Omega \\
U_s &= \frac{1}{2} \int_\Omega D^s \left[ \left( \frac{\partial W}{\partial x} + \Psi_x \right)^2 + \left( \frac{\partial W}{\partial y} + \Psi_y \right)^2 \right] d\Omega \\
T_{\max} &= \frac{1}{2} \omega^2 \int_\Omega \left( \rho h W^2 + I_2 (\Psi_x^2 + \Psi_y^2) \right) d\Omega
\end{align}
總能量泛函 $\Pi = U_b + U_s - T_{\max}$。

\subsection{Ritz 法離散}
將 $W,\Psi_x,\Psi_y$ 展開為基函數線性組合：
\begin{align}
W &= \sum_{j=1}^N W_j \Phi_j^w,\quad
\Psi_x = \sum_{j=1}^N X_j \Phi_j^x,\quad
\Psi_y = \sum_{j=1}^N Y_j \Phi_j^y
\end{align}
對泛函取變分 $\partial \Pi / \partial W_i = 0$, $\partial \Pi / \partial X_i = 0$, $\partial \Pi / \partial Y_i = 0$，整理成廣義特徵值問題：
\begin{equation}
\begin{bmatrix}
\mathbf{K}_{ww} & \mathbf{K}_{w\psi_x} & \mathbf{K}_{w\psi_y} \\
\mathbf{K}_{\psi_x w} & \mathbf{K}_{\psi_x\psi_x} & \mathbf{K}_{\psi_x\psi_y} \\
\mathbf{K}_{\psi_y w} & \mathbf{K}_{\psi_y\psi_x} & \mathbf{K}_{\psi_y\psi_y}
\end{bmatrix}
\begin{Bmatrix} \{W\} \\ \{X\} \\ \{Y\} \end{Bmatrix}
= \omega^2
\begin{bmatrix}
\mathbf{M}_{ww} & 0 & 0 \\
0 & \mathbf{M}_{\psi_x\psi_x} & 0 \\
0 & 0 & \mathbf{M}_{\psi_y\psi_y}
\end{bmatrix}
\begin{Bmatrix} \{W\} \\ \{X\} \\ \{Y\} \end{Bmatrix}
\end{equation}
各子矩陣由基函數積分給出（詳見推導）。此形式可用有限元素法或無網格法數值求解。